

\begin{frame}
	\frametitle{Scripts de règles}

\end{frame}

\begin{frame}
	\frametitle{Extension aux scripts}
	\begin{block}{Idée}
		Dans cette slide, remettre en contexte avec la question des opérations
		complexes. %
		Les modeleurs paramétriques existants proposent des opérations basées sur du
		code et comprennent nativement des structures de contrôle (séquence,
		répétition (itération), alternative). %
		Notre troisième contribution est l'extension de notre mécanisme de rejeu pour le support de scripts. %
		Serait-ce une bonne idée de représenter ces structures à l'aide de schémas ?
	\end{block}
	\begin{itemize}
		\item Extension de notre mécanisme de rejeu aux scripts,
		\item trois structures de contrôle (séquence, itération, alternative)
	\end{itemize}
\end{frame}

\begin{frame}
	\frametitle{Structure de séquence}
	\begin{itemize}
		\item Problématique de la représentation d'un script boîte fermée vs boîte ouverte,
		\item extension pour les spécifications paramétriques,
		\item extension pour les DAG
	\end{itemize}
\end{frame}

\begin{frame}
	\frametitle{Structure de boucle d'itération de sous-orbites}
	\begin{itemize}
		\item Problématique de l'itération de sous-orbite avec conditionnement sur
		      l'orbite itéré et ses sous-orbites d'itération,
		\item extension pour les spécifications paramétriques,
		\item extension pour les DAG
	\end{itemize}
\end{frame}

\begin{frame}
	\frametitle{Structure d'alternative de règles}
	\begin{itemize}
		\item Problématique de l'alternative avec des règles qui peuvent filtrer
		      différemment et appliquer des transformations différentes,
		\item extension pour les spécifications paramétriques,
		\item extension pour les DAG
	\end{itemize}
\end{frame}

\begin{frame}
	\frametitle{Appariement de règles alternatives}
	\begin{block}{Note}
		Je manque encore trop de recul, besoin d'en discuter. %
	\end{block}
	\begin{itemize}
		\item Définition de classe d'évènements,
		\item filtrage droit de la règle évaluée,
		\item filtrage gauche des deux règles,
		\item appariement à droite sur la base des origines qui peuvent être
		      strictement identiques ou des sur ou sous-orbites et à condition que
		      les orbites à apparier soient définies appartiennent à une même
		      classe d'évènements topologiques
	\end{itemize}
\end{frame}




%%% Local Variables:
%%% mode: LaTeX
%%% TeX-master: "../../main"
%%% End:
